%============================================================
% SISTEMATIKA PENULISAN BUKU KP (5 Bab)
% Di-input oleh chapters/bab1.tex via \input{\SistematikaFile}
%============================================================

\vspace{0.3cm}
\begin{description}[leftmargin=2.2cm, labelsep=0cm, itemsep=0.1cm]
  \item[{\makebox[2.2cm][l]{\textbf{BAB I}}}] \textbf{Pendahuluan}\\
    Bab ini memuat penjelasan mengenai latar belakang dilaksanakannya kerja
    praktik, perumusan masalah, batasan ruang lingkup, tujuan dan manfaat dari
    kerja praktik. Pada akhir bab juga disajikan sistematika penulisan untuk
    memberikan gambaran menyeluruh mengenai struktur laporan secara keseluruhan.

  \item[{\makebox[2.2cm][l]{\textbf{BAB II}}}] \textbf{Gambaran Perusahaan}\\
    Bab ini menjelaskan profil \NamaPerusahaan\ sebagai tempat pelaksanaan
    kerja praktik. Uraian mencakup struktur organisasi perusahaan, pembagian
    divisi kerja, hak dan wewenang tiap jabatan penting, lokasi perusahaan,
    serta penerapan etika profesi di lingkungan kerja.

  \item[{\makebox[2.2cm][l]{\textbf{BAB III}}}] \textbf{Landasan Teori}\\
    Bab ini berisi landasan teoritis yang menjadi acuan dalam pelaksanaan kerja
    praktik. Teori-teori yang dibahas meliputi teknologi dan framework yang
    digunakan dalam proyek. Penjelasan mencakup konsep dasar, kelebihan, dan
    implementasi umum dari masing-masing teknologi.

  \item[{\makebox[2.2cm][l]{\textbf{BAB IV}}}] \textbf{Rancangan dan Implementasi Sistem}\\
    Bab ini menjelaskan proses perancangan dan implementasi sistem selama kerja
    praktik. Pembahasan mencakup standar perusahaan dalam pengembangan perangkat
    lunak serta implementasi proyek yang dikerjakan, termasuk latar belakang,
    tujuan, deskripsi aplikasi, dan perancangan sistem.

  \item[{\makebox[2.2cm][l]{\textbf{BAB V}}}] \textbf{Penutup}\\
    Bab terakhir berisi kesimpulan dari kegiatan kerja praktik yang telah
    dilaksanakan serta hasil yang diperoleh. Selain itu, diberikan pula saran
    sebagai bahan masukan bagi perusahaan, institusi pendidikan, dan mahasiswa
    untuk pengembangan selanjutnya.
\end{description}
